\section{Responsible Research}
\label{sec:responsible_research}

\textbf{Reproducibility.} Every pre-training run is indexed by code commit,
split size, and seed. Each checkpoint stores the git revision that produced
it, and resuming a run across code versions triggers an explicit warning, so
no reported number mixes two versions of the training code. The data
pipeline is deterministic given its seed: the nested splits are prefixes of
one published shuffled index list, and training fixes all framework seeds
with deterministic cuDNN kernels. One precision exception is disclosed --- pre-training
matrix multiplications use TF32, which is reproducible on the same
hardware but not bit-identical to fp32 across machines; the linear-probe
evaluation runs in fp32 --- and Section~\ref{sec:limitations} discloses a
cross-revision provenance caveat for the largest splits. All hyperparameters appear in
Sections~\ref{sec:method} and~\ref{sec:setup}, including deviations from
published defaults. The code, split indices, evaluation scripts, and figure
scripts are kept under version control alongside the logged metrics of
every run, and one consumer GPU reproduces any single curve point in hours
(Section~\ref{sec:discussion}).

\textbf{Honest reporting.} Checkpoint selection never sees the downstream
dataset, result claims are scoped to best-within-budget readings, and the
known metric inflations and untested constants are collected in
Section~\ref{sec:limitations} with their direction stated. Interpretations
are labeled as interpretations, and unflattering outcomes --- a curve that
dips, splits whose validation had not plateaued, an inflation we cannot yet
quantify --- are reported as such.

\textbf{Data and ethics.} The study uses two long-established public
research datasets, Tiny-ImageNet (an ImageNet derivative) and CIFAR-10,
under their existing research terms. No new data is collected and no human
subjects are involved. The datasets inherit the documented limitations of
their sources, including category and imagery biases; nothing in this study
audits or corrects them, and the learned encoders are research artifacts for
measuring data efficiency, not models intended for deployment. The compute
footprint is deliberately small --- a single consumer GPU for hours rather
than GPU-years. Keeping the study reproducible at that scale is part of the
research setting itself: making the behavior of self-supervised methods
measurable with resources most groups have.
